\documentclass[12pt, a4paper]{article}

\usepackage{amsmath}
\usepackage{amssymb}
\usepackage{titlesec}
\usepackage{tcolorbox}
\usepackage{enumitem}
\usepackage[bottom=3cm]{geometry}
\usepackage[hidelinks]{hyperref}
\usepackage{forest}
\usepackage{pgfplots}
\tcbuselibrary{breakable}
\tcbuselibrary{skins}

\newtcolorbox{prob}[1]{colback=gray!5!white, colframe=gray!75!black, 
title=\textbf{Exercise #1}}

\newtcolorbox{sol}{
    breakable,
    colback=white,      % Background matches page
    colframe=white,     % Frame matches page (invisible)
    frame hidden,       % Hides the border line
    left=3mm, right=3mm,% MATCHES the Question box padding
    boxrule=0mm,        % No border width
    top=0mm, bottom=0mm,% Tight vertical spacing
    parbox=false,       % Allows paragraphs to break normally
    before upper={\textbf{Solution:}\par\medskip} % Automatically adds "Solution:" title
}


\begin{document}

\begin{prob}{5.8}
Is $C_2$ uniquely decodeable?
\end{prob}
\begin{sol}
Yes, because we can determine all the position where 101 is from where 0s are. Then, we can determine that the remaining 1s are coded from 1s.
The codewords can be uniquely decomposed and decoded.
\end{sol}
\bigskip


\begin{prob}{5.14}
\end{prob}
\begin{sol}
\end{sol}
\bigskip


\begin{prob}{5.16}
\end{prob}
\begin{sol}
\end{sol}
\bigskip


\begin{prob}{5.19}
Is the code $\{00, 11, 0101, 111, 1010, 100100, 0110\}$ uniquely decodeable?
\end{prob}
\begin{sol}
No, since '111111' can be decoded as '11/11/11' or '111/111'.
\end{sol}
\bigskip


\begin{prob}{5.20}
Is the ternary code $\{00, 012, 0110, 0112, 100, 201, 212, 22\}$ uniquely decodeable?
\end{prob}
\begin{sol}
This is a prefix code so yes, it is uniquely decodable.
\end{sol}
\bigskip

% template
\begin{prob}{}
\end{prob}
\begin{sol}
\end{sol}
\bigskip


\end{document}



